\documentclass[12pt,a4paper]{amsart}

\usepackage{amsmath,amssymb,amsthm}
\usepackage{mathtools}
\usepackage{enumitem}
\usepackage{hyperref}

\newtheorem{theorem}{Theorem}[section]
\newtheorem{lemma}[theorem]{Lemma}
\newtheorem{proposition}[theorem]{Proposition}
\newtheorem{corollary}[theorem]{Corollary}
\theoremstyle{definition}
\newtheorem{definition}[theorem]{Definition}
\newtheorem{assumption}[theorem]{Assumption}
\newtheorem{hypothesis}[theorem]{Hypothesis}
\theoremstyle{remark}
\newtheorem{remark}[theorem]{Remark}

\newcommand{\Rig}{\textup{[Rig]}}
\newcommand{\Cond}{\textup{[Cond]}}

\begin{document}

\title[A$_2$-Stability for PSWF Amplitudes]{%
  Proof Architecture: A$_2$-Stability for Prolate Spheroidal Wave Function Amplitudes\\
  in the Semiclassical Limit}

\author{[Author]}

\begin{abstract}
This paper establishes the proof architecture for the A$_2$-stability
of PSWF amplitudes in the semiclassical limit $h:=c^{-1}\to 0$
with $n\ge 0$ fixed.

\medskip
\noindent\textbf{Objective.}
Verify Hypothesis~\textup{hyp:rigidity} from Paper~XVI:
the PSWF amplitude satisfies Assumption~A2 uniformly as $h\to 0$.

\medskip
\noindent\textbf{Assumption~A2.}
(A1) $c_3(h)\ge\delta_0>0$;
(A2) $|\det(dt/d\beta)|\ge\delta_1>0$;
(A3) unique non-degenerate fold ($V'(x_0)\ne 0$; standard A$_2$ condition).

\medskip
\noindent\textbf{Semiclassical regime.}
$n\ge 0$ \emph{fixed}; $h=c^{-1}\to 0$.
No Plancherel--Rotach or double-scaling regime is invoked;
all asymptotics are in fixed quantum number scaling.

\medskip
\noindent\textbf{Three hard cores.}
\begin{enumerate}[\rm(C1)]
\item Spectral gap: $E_{n+1}(h)-E_n(h)\ge Ch$
  (Lemma~\ref{lem:anchor} + Theorem~\ref{thm:specgap}).
\item Smooth eigenbranch: $h\mapsto E_n(h)$ real-analytic
  (Theorem~\ref{thm:eigenbranch}).
\item Olver $C^1$-stability: $\inf_{K_\varepsilon}|\phi'(\cdot;h)|\ge\delta>0$
  (Lemmas~\ref{lem:C0-airy}--\ref{lem:tp-perturbation}).
\end{enumerate}

\medskip
\noindent\textbf{Status.}
The proof graph is structurally complete and acyclic.
Two technical write-ups remain open
(explicit computation of $T(\lambda_n)$ and
CFU Jacobian ellipticity; see Remark~\ref{rem:remaining}).
\end{abstract}

\maketitle
\tableofcontents

%% ============================================================
\section{Input from Paper~XVI}
\label{sec:input}
%% ============================================================

\begin{itemize}
\item \textbf{thm:CFU-geom} (global gluing theorem):
  normalized Airy coordinates $\Phi=\phi\theta-\tfrac13\theta^3$
  on each $K_\varepsilon^\pm$ via smooth canonical transformation;
  branches glued with Maslov index $+1$
  (Remark~\ref{rem:cfu-geom-nontrivial}).
\item \textbf{thm:CFU-symbol}: $I^{-1/4,0}\to S^{-1/2}_{1,0}(h)$, uniform.
\item \textbf{Assumption~A2}: see Abstract.
\item \textbf{hyp:rigidity}: verified here.
\end{itemize}

\begin{remark}[CFU-geom: gluing and phase fixation]
\label{rem:cfu-geom-nontrivial}
In 1D with a single non-degenerate fold (Hypothesis~\ref{hyp:a2-fold}):
(i) Maslov index $+1$ across $x_0(h)$ (\cite{Zworski2012} Ch.~12). \Rig.
(ii) No phase mismatch: $\phi$ smooth with fixed sign on each $K_\varepsilon^\pm$. \Rig.
(iii) \emph{Phase fixation.}
  The local Airy normal form $\Phi=\phi\theta-\tfrac13\theta^3$
  determines $\phi$ up to sign on each $K_\varepsilon^\pm$.
  The global constant phase freedom is fixed by the
  $L^2$-normalization of the eigenfunction: $\|\psi_n\|_{L^2}=1$.
  Hence the CFU representation is
  \emph{well-defined up to a global constant fixed by eigenfunction normalization},
  not uniquely canonical in the abstract sense.
  This suffices for (A2): $|\phi'|$ is sign-independent. \Rig.
(iv) Multiple turning points: each A$_2$-fold treated separately. \Rig.
\end{remark}

%% ============================================================
\section{PSWF Setup and Standing Assumptions}
\label{sec:pswf}
%% ============================================================

\begin{definition}[PSWF, classical and semiclassical momentum]
\label{def:pswf-p0}
$\psi_n$ satisfy $(h^2\mathcal{L}_{h^{-1}})\psi_n=E_n(h)\psi_n$,
$V(x)=x^2/(1-x^2)$, $E_n(h):=\chi_n h^2$.
$n\ge 0$ \emph{fixed}; $h\to 0$.

\noindent\textbf{Limiting energy:}
$\lambda_n:=\lim_{h\to 0}E_n(h)$ (Theorem~\ref{thm:eigenbranch}).

\noindent\textbf{Classical momentum:}
\begin{equation}\label{eq:p0-def}
  p_0(x):=\sqrt{\lambda_n-V(x)},\quad x\in\{\lambda_n>V(x)\}.
\end{equation}

\noindent\textbf{Semiclassical momentum:}
$p(x;h):=\sqrt{E_n(h)-V(x)}$ for $x\in K_\varepsilon$, $h\le h_0$.
\end{definition}

\begin{hypothesis}[Non-degenerate fold: standard A$_2$ condition]
\label{hyp:a2-fold}
The limiting turning point $x_0^{(0)}$ satisfies:
\begin{enumerate}[{\rm(F1)}]
\item $V'(x_0^{(0)})\ne 0$.
  \emph{Verification:} $V'(x)=2x/(1-x^2)^2\ne 0$
  for $x\in(0,1-\delta_{\rm bdy})$. \Rig.
\end{enumerate}
(F1) is the standard A$_2$ fold condition in the Hamilton-Jacobi / CFU sense
\cite{Zworski2012}: non-degenerate turning point $\Leftrightarrow$ $V'(x_0)\ne 0$.
The second-order structure follows from Lemma~\ref{lem:vpp} below (not a hypothesis).
\end{hypothesis}

\begin{lemma}[Second-order non-degeneracy of $V$]
\label{lem:vpp}
$V''(x)=2(1+3x^2)/(1-x^2)^3>0$ for all $x\in(-1,1)$. \Rig.
Consequences:
\begin{enumerate}[{\rm(i)}]
\item Non-degenerate quadratic term at $x_0^{(0)}$:
  $V''(x_0^{(0)})>0$. \Rig.
\item Airy normal form stability: the canonical coordinate change
  of thm:CFU-geom is a local $C^\infty$ diffeomorphism
  (Morse lemma stability; \cite{Olver1974} Ch.~10). \Cond.
\item Ensures $|I|,|\phi|\ge c>0$ on $K_\varepsilon$,
  used in Lemmas~\ref{lem:nondeg} and~\ref{lem:deriv-commuting}. \Cond.
\end{enumerate}
\emph{Not a hypothesis; follows from the explicit formula for $V$
and Assumption~\ref{ass:boundary-away}.}
\end{lemma}

\begin{remark}[Shape of $V$]
\label{rem:V-shape}
\begin{enumerate}[{\rm(i)}]
\item $V(0)=0$, $V\to+\infty$ as $x\to\pm 1$, $V(-x)=V(x)$.
\item $V'>0$ on $(0,1)$; $<0$ on $(-1,0)$. \Rig.
\item Unique turning point per half-interval
  for $\lambda_n<V(1-\delta_{\rm bdy})$. \Rig.
\item Multiple turning points: each a separate A$_2$-fold. \Rig.
\end{enumerate}
\end{remark}

\begin{remark}[Geometric positivity]
\label{rem:geo-positivity}
$c_\varepsilon^{\rm geo}:=\inf_{K_\varepsilon}(\lambda_n-V)>0$. \Rig.
\textbf{All constants $h$-independent.}
\end{remark}

\begin{assumption}[Boundary-away cutoff]
\label{ass:boundary-away}
$\mathrm{supp}\,\chi\subset[-1+\delta_{\rm bdy},1-\delta_{\rm bdy}]$,
indep.\ of $h$; $V,V',V''\in L^\infty$;
$\lambda_n<V(1-\delta_{\rm bdy})$.
\end{assumption}

%% ============================================================
\section{Hard Core (C2): Smooth Eigenbranch}
\label{sec:eigenbranch}
%% ============================================================

\begin{theorem}[Smooth eigenbranch, \Cond]
\label{thm:eigenbranch}
Assume:
\begin{enumerate}[{\rm(E1)}]
\item \textit{(Fixed domain.)} $\mathcal{D}=H^2(-1,1)\cap H^1_0(-1,1)$,
  $h$-independent. \Rig\ ($V\in L^\infty$; boundary cond.\ independent of $h$).
\item \textit{(Kato--Rellich.)}
  Write $h^2\mathcal{L}_{h^{-1}}=\mathcal{L}_0+hA_1+h^2A_2$,
  $\mathcal{L}_0=-d^2/dx^2$.
  $A_1,A_2$ are multiplication operators with $L^\infty$-coefficients
  (Assumption~\ref{ass:boundary-away}),
  hence $\mathcal{L}_0$-relatively bounded with relative bound $a=0$. \Rig.
\item \textit{(Local resolvent analyticity.)}
  Fix $h_0>0$ and $z_0$ with
  $\mathrm{dist}(z_0,\sigma(h_0^2\mathcal{L}_{h_0^{-1}}))\ge\delta_z>0$
  (Lemma~\ref{lem:anchor}(Anc3) provides initial gap).
  Resolvent identity centered at the \emph{perturbed} operator:
  \begin{align*}
    (h^2\mathcal{L}_{h^{-1}}-z)^{-1}
    &= (h_0^2\mathcal{L}_{h_0^{-1}}-z)^{-1}\\
    &\quad\cdot\bigl[I+(h^2\mathcal{L}_{h^{-1}}-h_0^2\mathcal{L}_{h_0^{-1}})
      (h_0^2\mathcal{L}_{h_0^{-1}}-z)^{-1}\bigr]^{-1}.
  \end{align*}
  The perturbation is $\mathcal{O}(|h-h_0|)$ in $\mathcal{B}(L^2)$ by (E2);
  Neumann series converges in $\mathcal{B}(L^2)$ for $|h-h_0|$ small,
  analytic in $h$. \Cond.
\item \textit{(Simple eigenvalue.)} $E_n(h_0)$ simple: SL oscillation. \Rig.
\end{enumerate}
Hence $h^2\mathcal{L}_{h^{-1}}$ is a Kato analytic family of type~(A)
(\cite{Kato1966} Ch.~VII.\S2);
$h\mapsto E_n(h)$ is real-analytic with $E_n(h)\to\lambda_n$;
$x_0(h)$ real-analytic by IFT + (F1).
\end{theorem}

\begin{remark}[Local resolvent vs.\ free resolvent]
Centering the Neumann series at $h_0^2\mathcal{L}_{h_0^{-1}}$
avoids the hidden spectral gap between the free operator $\mathcal{L}_0$
and the perturbed family near $E_n(h)$.
The initial gap (Anc3) ensures $\delta_z>0$;
relative boundedness (E2) gives $\mathcal{O}(|h-h_0|)$ operator-norm control,
closing (E3) without reference to the free resolvent.
\end{remark}

%% ============================================================
\section{Hard Core (C1): Spectral Gap}
\label{sec:specgap}
%% ============================================================

\begin{lemma}[Anchor: variational + Agmon, \Cond]
\label{lem:anchor}
$n$ fixed, $h_0>0$ small:
\begin{enumerate}[{\rm(Anc1)}]
\item \textit{(Variational lower bound.)}
  By harmonic approximation of $V$ near the global minimum $x=0$
  (where $V(x)=x^2+\mathcal{O}(x^4)$) and the
  semiclassical harmonic oscillator spectrum:
  $E_n(h)=\mathcal{O}(h)$ with explicit lower bound
  $E_n(h)\ge c_n h$ for some $c_n>0$ depending only on $n$
  (\cite{Zworski2012} Ch.~6). \Cond.
  In particular $\{E_n(h)\}_{h\le h_0}$ lies in a compact interval $(0,C_n h_0]$.
\item \textit{(Agmon decay.)}
  $\|e^{d(\cdot,K_\varepsilon)/h}\psi_n\|_{L^2}=\mathcal{O}(1)$
  where $d$ is the Agmon distance for $V-\lambda_n$
  (\cite{Zworski2012} Thm.~6.3).
  Hence $\psi_n$ is supported in $K_\varepsilon$ up to $\mathcal{O}(e^{-c/h})$. \Cond.
\item \textit{(Initial gap at $h_0$.)}
  $E_{n+1}(h_0)-E_n(h_0)\ge\delta_{h_0}>0$
  by SL oscillation / variational separation. \Rig.
\end{enumerate}
\emph{Role:} (Anc1) initializes the WKB energy range;
(Anc2) justifies single-well WKB and controls tunneling;
(Anc3) provides $\delta_z>0$ for Theorem~\ref{thm:eigenbranch}(E3),
breaking the apparent circularity
WKB $\Rightarrow$ gap $\Rightarrow$ WKB.
\end{lemma}

\begin{lemma}[Uniform WKB quantization validity, \Cond]
\label{lem:wkb-validity}
$n$ fixed, $h\to 0$.
\emph{No Plancherel--Rotach regime; all asymptotics in fixed-$n$ scaling.}
\begin{enumerate}[{\rm(W1)}]
\item \textit{(Single well.)} $V$ has unique minimum at $x=0$;
  $E_n(h)\in(0,V(1-\delta_{\rm bdy}))$ by (Anc1);
  classically allowed interval $(x_-(h),x_+(h))$ connected
  and bounded away from $\pm 1$. \Rig.
\item \textit{(No tunneling.)} Tunneling $\mathcal{O}(e^{-c/h})$
  by (Anc2); does not affect $\mathcal{O}(h^2)$ in BS. \Cond.
\item \textit{(Olver uniformity.)} For $n$ fixed as $h\to 0$:
  $S(E_n(h)):=\int_{x_-}^{x_+}\sqrt{E_n-V}\,dx
  =(n+\tfrac12)\pi h+\mathcal{O}(h^2)$
  uniformly
  (Olver \cite{Olver1974} Ch.~10--11). \Cond.
\end{enumerate}
\end{lemma}

\begin{lemma}[Bohr--Sommerfeld inversion, \Cond]
\label{lem:bs-inversion}
Under (H2)--(H3) and Lemma~\ref{lem:wkb-validity}(W3):
\begin{enumerate}[{\rm(i)}]
\item $S'(E)=T(E)/2>0$: strict monotonicity near $\lambda_n$.
\item $S$ is a local $C^\infty$ diffeomorphism near $\lambda_n$.
\item \textit{(MVT with uniform denominator.)} For
  $\xi\in(E_n,E_{n+1})\to\lambda_n$ as $h\to 0$:
  $S'(\xi)\ge T(\lambda_n)/4>0$ for $h\le h_0$; hence
  \[
    E_{n+1}-E_n
    =\frac{S(E_{n+1})-S(E_n)}{S'(\xi)}
    =\frac{\pi h+\mathcal{O}(h^2)}{T(\lambda_n)/2+\mathcal{O}(h)}
    =\frac{2\pi h}{T(\lambda_n)}+\mathcal{O}(h^2).
  \]
  $T(\lambda_n)$ depends only on $\lambda_n$, not on $h$;
  hence $C=\pi/T(\lambda_n)>0$ is an $h$-independent constant.
\end{enumerate}
\end{lemma}

\begin{theorem}[Spectral gap, \Cond]
\label{thm:specgap}
Assume:
\begin{enumerate}[{\rm(H1)}]
\item 1D integrability. \Rig.
\item $0<T(\lambda_n)<\infty$: simple zeros of $V-\lambda_n$. \Cond.
\item $\partial_E T(\lambda_n)
  =-\tfrac12\oint(\lambda_n-V)^{-3/2}\,dx\ne 0$. \Cond.
\end{enumerate}
By Lemmas~\ref{lem:anchor},~\ref{lem:wkb-validity},~\ref{lem:bs-inversion}:
\begin{equation}\label{eq:specgap}
  E_{n+1}(h)-E_n(h)
  =\frac{2\pi h}{T(\lambda_n)}+\mathcal{O}(h^2)\ge Ch,
  \quad C=\frac{\pi}{T(\lambda_n)}>0.
\end{equation}
\end{theorem}

%% ============================================================
\section{Branch Selection and Microlocal Bridge}
\label{sec:branch}
%% ============================================================

\begin{proposition}[Branch selection, \Cond]
\label{prop:branch-selection}
\textit{(i)} $|E_n-\lambda_n|=\mathcal{O}(h)$;
$p\ge\delta_\xi:=\sqrt{c_\varepsilon^{\rm geo}/2}>0$ on $K_\varepsilon$.

\textit{(ii)} $\mathrm{WF}_h(\psi_n)\subset\Lambda_h:=\{\xi^2+V=E_n\}$
(\cite{Zworski2012} Thm.~12.1;
ellipticity of $h^2\mathcal{L}_{h^{-1}}-E_n$ outside $\Lambda_h$);
spectral gap (Thm.~\ref{thm:specgap}) excludes mode mixing;
$\chi\psi_n=\Pi^+(\chi\psi_n)+\mathcal{O}(h^\infty)$,
$\Pi^+(\chi\psi_n)\in I^{-1/4,0}(\Lambda_h^+)$.
\end{proposition}

\begin{remark}[Microlocal bridge: PDE $\to$ FIO]
\label{rem:micro-bridge}
$\mathrm{WF}_h(\psi_n)\subset\Lambda_h$ closes the PDE$\to$FIO step.
Stability under $h\to 0$: smooth eigenbranch (Thm.~\ref{thm:eigenbranch})
gives $\Lambda_h\to\Lambda_0$ smoothly;
spectral gap (Thm.~\ref{thm:specgap}) prevents mixing with $\psi_{n\pm 1}$.
Without this, the FIO structure of $\psi_n$ would be asserted, not derived.
\end{remark}

%% ============================================================
\section{CFU Construction on $K_\varepsilon^\pm$}
\label{sec:cfu-components}
%% ============================================================

\begin{remark}[CFU on each component]
\label{rem:cfu-components}
$K_\varepsilon^\pm:=K_\varepsilon\cap\{\pm(x-x_0(h))>0\}$.
On $K_\varepsilon^+$: $I>0$, $\phi=(\tfrac32 I)^{2/3}$;
on $K_\varepsilon^-$: $I<0$, $\phi=-(\tfrac32|I|)^{2/3}$.
CFU \cite{CFU1957} applied separately;
glued via thm:CFU-geom (Remark~\ref{rem:cfu-geom-nontrivial}).
\end{remark}

%% ============================================================
\section{Hard Core (C3): Olver $C^1$-Stability}
\label{sec:olver-C1}
%% ============================================================

\begin{lemma}[Uniform ellipticity, \Cond]
\label{lem:uniform-ellipticity}
$p\ge\delta_\xi$, $p_0\ge c_\varepsilon:=\sqrt{c_\varepsilon^{\rm geo}}$,
$p+p_0\ge\delta_\xi+c_\varepsilon>0$, uniform in $h$.
\end{lemma}

\begin{lemma}[Energy comparison, \Cond]
\label{lem:energy-comparison}
$|p-p_0|\le|E_n-\lambda_n|/(\delta_\xi+c_\varepsilon)=:A_\varepsilon=\mathcal{O}(h)$.
\end{lemma}

\begin{lemma}[Integral stability, \Cond]
\label{lem:integral-stability}
$|I-I_0|\le\|p\|_{C^0}|x_0(h)-x_0^{(0)}|+|K_\varepsilon|A_\varepsilon
=:\tilde A_\varepsilon=\mathcal{O}(h)$.
\end{lemma}

\begin{lemma}[Non-degeneracy, \Cond]
\label{lem:nondeg}
$|I|,|I_0|\ge c_\varepsilon':=\delta_\xi\varepsilon/2>0$;
$|\phi|\ge c_\varepsilon''$; $|\phi_0|\ge d_\varepsilon>0$.
\emph{Uses Lemma~\ref{lem:vpp}(iii): $V''>0$ ensures
non-degenerate quadratic structure of $V$ near $x_0^{(0)}$,
which prevents $I$ from vanishing on $K_\varepsilon$.}
\end{lemma}

\begin{lemma}[$C^0$-Airy stability, \Cond]
\label{lem:C0-airy}
$\|\phi-\phi_0\|_{C^0}\le L_\varepsilon\tilde A_\varepsilon=\mathcal{O}(h)$;
$L_\varepsilon=\tfrac23(\tfrac32)^{2/3}(c_\varepsilon')^{-1/3}$
(Lipschitz of $t\mapsto(\tfrac32 t)^{2/3}$ on $[c_\varepsilon',\infty)$;
stabilized by Lemma~\ref{lem:vpp}(ii)).
\end{lemma}

\begin{lemma}[Derivative commuting, \Cond]
\label{lem:deriv-commuting}
On $K_\varepsilon^\pm$:
$\phi'=p/\phi$, $\phi_0'=p_0/\phi_0$ (algebraic formulas).
$\|\phi'-\phi_0'\|_{C^0}=\mathcal{O}(h)$ by quotient estimate.
\emph{Uses Lemma~\ref{lem:vpp}(iii): $\phi\ne 0$ on $K_\varepsilon$.}
\end{lemma}

\begin{lemma}[Turning-point perturbation, \Cond]
\label{lem:tp-perturbation}
$|x_0(h)-x_0^{(0)}|=(E_n-\lambda_n)/V'(x_0^{(0)})+\mathcal{O}(h^2)=\mathcal{O}(h)$
(IFT + Theorem~\ref{thm:eigenbranch} + (F1)).
\end{lemma}

\begin{corollary}[$C^1$-stability]
\label{cor:phi-C1}
$\inf_{K_\varepsilon}|\phi'(\cdot;h)|\ge\tfrac12\inf|\phi_0'|>0$;
$\phi_0'=p_0/\phi_0\ge c_\varepsilon/d_\varepsilon>0$.
\end{corollary}

%% ============================================================
\section{CFU Jacobian in 1D}
\label{sec:cfu-jacobian}
%% ============================================================

\begin{lemma}[CFU Jacobian: symplectic control, \Cond]
\label{lem:cfu-jacobian-1d}
In Airy normal form coordinates $\Phi=\phi\theta-\tfrac13\theta^3$
(phase fixed up to global constant by $L^2$-normalization;
Remark~\ref{rem:cfu-geom-nontrivial}(iii)):
\begin{enumerate}[{\rm(i)}]
\item $\partial^2_{x\theta}\Phi=\phi'$. \Rig.
\item \textit{(Symplectic factor in normalized coordinates.)}
  In normalized Airy coordinates, the canonical transformation
  $(x,\xi)\mapsto(x,\theta)$ of thm:CFU-geom satisfies
  $|\partial(x,\xi)/\partial(x,\theta)|=|\phi'|$
  (\cite{Olver1974} Ch.~10; \cite{Zworski2012} Ch.~12).
  Any remaining smooth amplitude factor arising from the
  canonical transformation is nonvanishing and bounded away from $0$;
  it is absorbed into the symbol class $S^{-1/2}_{1,0}(h)$
  by thm:CFU-symbol and does not affect the lower bound (A2). \Cond.
\item Hence $|\det(dt/d\beta)|=|\phi'|\cdot(1+\mathcal{O}(h))\ge\delta_1>0$
  (Corollary~\ref{cor:phi-C1}). \Cond.
\end{enumerate}
\end{lemma}

%% ============================================================
\section{Derivation of Assumption~A2}
\label{sec:a2-derivation}
%% ============================================================

\begin{lemma}[Microlocal Structure, \Cond]
\label{lem:olver-fio-lift-micro}
(i) $\Pi^+(\chi\psi_n)\in I^{-1/4,0}(\Lambda_h^+)$
(Remark~\ref{rem:micro-bridge}).
(ii) Unique A$_2$-fold; $c_3(h)\ge\delta_0>0$
(Hypothesis~\ref{hyp:a2-fold}(F1) + Lemma~\ref{lem:vpp}).
\end{lemma}

\begin{proposition}[Assumption~A2]
\label{prop:a2-from-lemmas}
(A1): Lemma~\ref{lem:olver-fio-lift-micro}(ii).
(A2): Lemma~\ref{lem:cfu-jacobian-1d}(iii) + Corollary~\ref{cor:phi-C1}.
(A3): Hypothesis~\ref{hyp:a2-fold}(F1) + Remark~\ref{rem:V-shape}(ii)--(iii).
\end{proposition}

%% ============================================================
\section{Status Table}
\label{sec:status}
%% ============================================================

\begin{center}
\small
\begin{tabular}{lllp{3.5cm}}
\hline
Statement & Core & Status & Source \\
\hline
(F1) $V'(x_0)\ne 0$ (A$_2$ standard) & --- & \Rig & explicit \\
$V''>0$ (auxiliary; not hypothesis) & --- & \Rig & Lem.~\ref{lem:vpp} \\
Airy normal form stable (Morse) & --- & \Cond & Lem.~\ref{lem:vpp}(ii) \\
$V$ monotone per half-interval & --- & \Rig & $V'$ formula \\
Unique TP; $\lambda_n<V(1-\delta)$ & --- & \Rig & Rem.~\ref{rem:V-shape} \\
$c_\varepsilon^{\rm geo}>0$ & --- & \Rig & Rem.~\ref{rem:geo-positivity} \\
Phase fixed by $L^2$-normalization & --- & \Rig & Rem.~\ref{rem:cfu-geom-nontrivial}(iii) \\
$|\phi'|$ sign-independent & --- & \Rig & Lem.~\ref{lem:cfu-jacobian-1d}(i) \\
CFU on $K_\varepsilon^\pm$ separately & --- & \Rig & Rem.~\ref{rem:cfu-components} \\
$p_0=\sqrt{\lambda_n-V}$ defined & --- & \Rig & Def.~\ref{def:pswf-p0} \\
$n$ fixed, $h\to 0$ regime & --- & \Rig & Def.~\ref{def:pswf-p0} \\
(E1) fixed domain $H^2\cap H^1_0$ & (C2) & \Rig & SL structure \\
(E2) $A_j$ multiplication, $a=0$ & (C2) & \Rig & $V\in L^\infty$ \\
(E3) local resolvent analytic & (C2) & \Cond & Thm.~\ref{thm:eigenbranch}(E3) \\
(E4) simple eigenvalue & (C2) & \Rig & SL oscillation \\
$E_n(h)$ real-analytic & (C2) & \Cond & Kato type~(A) \\
$x_0(h)$ real-analytic & (C2) & \Cond & IFT + (F1) \\
(Anc1) $E_n=\mathcal{O}(h)$; lower bound & (C1) & \Cond & Lem.~\ref{lem:anchor} \\
(Anc2) Agmon decay & (C1) & \Cond & Lem.~\ref{lem:anchor} \\
(Anc3) initial gap at $h_0$ & (C1) & \Rig & SL oscillation \\
(W1) single well; interval away & (C1) & \Rig & Lem.~\ref{lem:wkb-validity} \\
(W2) tunneling $e^{-c/h}$ & (C1) & \Cond & Agmon \\
(W3) Olver; BS formula ($n$ fixed) & (C1) & \Cond & Lem.~\ref{lem:wkb-validity} \\
(H2) $0<T(\lambda_n)<\infty$ & (C1) & \Cond & simple zeros \\
(H3) $\partial_E T\ne 0$ & (C1) & \Cond & explicit integral \\
$T(\lambda_n)$ $h$-independent & (C1) & \Rig & $\lambda_n=\lim E_n$ only \\
BS inversion; MVT uniform & (C1) & \Cond & Lem.~\ref{lem:bs-inversion} \\
Spectral gap $\ge Ch$ & (C1) & \Cond & Thm.~\ref{thm:specgap} \\
Microlocal bridge $\psi_n\in I^{-1/4,0}$ & --- & \Cond & Rem.~\ref{rem:micro-bridge} \\
$p\ge\delta_\xi$; branch separation & (C1)+(C2) & \Cond & Prop.~\ref{prop:branch-selection} \\
$p+p_0\ge c>0$ & (C3) & \Cond & Lem.~\ref{lem:uniform-ellipticity} \\
$A_\varepsilon,\tilde A_\varepsilon=\mathcal{O}(h)$ & (C3) & \Cond & Lems.~\ref{lem:energy-comparison},~\ref{lem:integral-stability} \\
$|I|,|\phi|\ge c>0$ (via $V''>0$) & (C3) & \Cond & Lem.~\ref{lem:nondeg} \\
$\|\phi-\phi_0\|_{C^0}=\mathcal{O}(h)$ & (C3) & \Cond & Lem.~\ref{lem:C0-airy} \\
Deriv.\ commuting algebraic & (C3) & \Cond & Lem.~\ref{lem:deriv-commuting} \\
$\inf|\phi'|\ge\delta>0$ & (C3) & \Cond & Cor.~\ref{cor:phi-C1} \\
Sympl.\ factor absorbed in symbol & (C3) & \Cond & Lem.~\ref{lem:cfu-jacobian-1d}(ii) \\
$|\det|\ge\delta_1>0$ & (C3) & \Cond & Lem.~\ref{lem:cfu-jacobian-1d}(iii) \\
(A1) $c_3\ge\delta_0$ & --- & \Cond & Lem.~\ref{lem:olver-fio-lift-micro} \\
(A2) $|\det|\ge\delta_1$ & (C3) & \Cond & Prop.~\ref{prop:a2-from-lemmas} \\
(A3) unique non-deg.\ fold & --- & \Rig & (F1) \\
\textup{hyp:rigidity} & all & \Cond & Thm.~\ref{thm:xvii-main} \\
\hline
\end{tabular}
\end{center}

\begin{remark}[Remaining tasks]
\label{rem:remaining}
\begin{enumerate}
\item (C1): explicit computation of $T(\lambda_n)$ and $\partial_E T(\lambda_n)$
  for $V=x^2/(1-x^2)$;
  turning points $x_\pm=\pm\sqrt{\lambda_n/(1+\lambda_n)}$ (algebraic);
  integral elliptic and explicit.
\item (C3): full write-up of CFU Jacobian ellipticity
  (Lem.~\ref{lem:cfu-jacobian-1d}(ii)): explicit separation of
  geometric Jacobian, Airy rescaling, and FIO amplitude correction;
  verification against \cite{Olver1974} Ch.~10 and \cite{Zworski2012} Ch.~12.
\item (C2): close local resolvent (E3) using (Anc3) as initialization.
\item (C3): full write-up of Lemmas~\ref{lem:C0-airy}--\ref{lem:tp-perturbation}.
\end{enumerate}
\end{remark}

%% ============================================================
\section{Target Theorem}
\label{sec:target}
%% ============================================================

\begin{theorem}[PSWF A$_2$-stability]
\label{thm:xvii-main}
Under Assumption~\ref{ass:boundary-away},
Hypothesis~\ref{hyp:a2-fold},
and hard cores (C1)--(C3),
the PSWF amplitude satisfies Assumption~A2
uniformly on $K_\varepsilon$ as $h\to 0$ ($n$ fixed).
\end{theorem}

\begin{remark}[Dependency chain]
\[
  \underbrace{\text{Anc}}_{\text{Lem.}~\ref{lem:anchor}}
  \Rightarrow
  \underbrace{\text{(C2)}}_{\text{Thm.}~\ref{thm:eigenbranch}}
  +\underbrace{\text{(C1)}}_{\text{Thm.}~\ref{thm:specgap}}
  +\underbrace{\text{Lem.}~\ref{lem:vpp}}_{V''>0}
  \;\Rightarrow\;
  \underbrace{\text{Prop.}~\ref{prop:branch-selection}+\text{Rem.}~\ref{rem:micro-bridge}}_{\psi_n\in I^{-1/4,0}}
  \;\Rightarrow\;
  \underbrace{\text{(C3)}}_{\phi\text{-chain}}
  \;\Rightarrow\;
  \text{Lem.}~\ref{lem:cfu-jacobian-1d}
  \;\Rightarrow\;
  \text{Thm.}~\ref{thm:xvii-main}.
\]
\end{remark}

\begin{thebibliography}{9}
\bibitem{PaperXVI}[Author], \textit{Paper~XVI}, 2026.
\bibitem{Olver1974}F.~W.~J.~Olver, \textit{Asymptotics and Special Functions}, Academic Press, 1974.
\bibitem{Zworski2012}M.~Zworski, \textit{Semiclassical Analysis}, AMS, 2012.
\bibitem{Hormander1983}L.~H\"ormander, \textit{The Analysis of Linear PDE~I}, Springer, 1983.
\bibitem{CFU1957}C.~Chester, B.~Friedman, F.~Ursell, Proc.\ Cambridge Philos.\ Soc.\ \textbf{53} (1957), 599--611.
\bibitem{MelroseU1979}R.~Melrose, G.~Uhlmann, Comm.\ Pure Appl.\ Math.\ \textbf{32} (1979), 483--531.
\bibitem{Slepian1961}D.~Slepian, H.~O.~Pollak, Bell Syst.\ Tech.\ J.\ \textbf{40} (1961), 43--63.
\bibitem{Kato1966}T.~Kato, \textit{Perturbation Theory for Linear Operators}, Springer, 1966.
\end{thebibliography}

\end{document}
